\section{Threats to Validity}

\begin{itemize}
  \item \textbf{Chronology.} RQ2 and RQ3 were posed after the RQ1 results; the initial RQ2 analysis was post hoc. Later protocols were written in advance but are internal.
  \item \textbf{Simulation.} RQ1 and the preregistered test are decision simulations designed by us; their outcomes are close to by-construction and yield no live-agent effect estimate. A small diagnostic with a real agent (Qwen2.5 7B) did not meet the condition for estimating an effect.
  \item \textbf{Independent corroborators are assumed truthful} apart from leakage. VP-Control reports that on FinQA small verifiers misread evidence and the source effect reverses; an independent source helps only if the verifier is competent.
  \item \textbf{Scope of the theorem.} It does not cover correlated labels, a wrong primary label, or multiple corroborators, and the proof has not been checked independently.
  \item \textbf{Survey labels.} They are AI-coded and the second coder also used AI. Agreement between AI coders shows that labels are reproducible, not that they are correct. The only human checks are the author's adjudication and the 10 evidence citations. A planned blind coding of 15 servers by the author was abandoned because AI assistance had been used to search the repositories; it is not reported as validation. The five survey-2 interpretation rules were set by coder A after reading code but before computing metrics.
  \item \textbf{Construct of sharing.} The operator level is a lower bound on true dependence; the non-empty-intersection rule overstates dependence for multi-source aggregators; the data-origin level uses only in-repository evidence.
  \item \textbf{Sample.} Popular GitHub MCP repositories only; pairs are not independent because each server appears in several pairs; the denominators of L1 recall are small (8 and 27); $\rho$ by type differs widely between samples.
  \item \textbf{Projections are not measurements of agents.}
\end{itemize}
